\subsection{Introduction to Factoring Polynomials} 
When we consider factoring polynomials, we will consider polynomials with \emph{integer} coefficients.
\begin{definition}
\begin{itemize*}
\item A \term{factor} of an integer $n$ is a number $k$ such that $n=kd$ for some integer $d$.
\item A \term{factor} of a monomial $M$ is a monomial $A$ such that $M=AB$ for some monomial $B$.
\item A \term{factor} of a polynomial $P$ is a polynomial $Q$ such that $P=QS$ for some polynomial $S$.
\end{itemize*}
\end{definition}
Note that these definitions are very similar!
\begin{example}
\item As an integer, $5$ is a factor of $30$ because $30=5\cdot\makeblanksmall$.
\item As a monomial, $3x^2$ is a factor of $9x^3y$ because $9x^3y=(3x^2)(\makeblanksmall)$.
\item As a polynomial, $4x+3$ is a factor of $4x^2+3x$ because $4x^2+3x=(\makeblanksmall)(4x+3)$. 
\end{example}
\noindent When we factor (verb) a polynomial, our goal is to rewrite it in a factored form.
\begin{definition}
An expression is in ``factored form'' if it is written as a product of simpler expressions, that is, if the last step to evaluate it would be \makeblank[1.5in].
\end{definition}
\begin{example}For each expression below, determine whether it is in a factored form.
\begin{enumerate}[a.]
\item\makeblankfill{$3x^2y$}
\item\makeblankfill{$(2x-3)(5x+2)$}
\item\makeblankfill{$3x(x+5)-7(x+5)$}
\item\makeblankfill{$4x^2+6x-3$}
\item\makeblankfill{$x^2(3x^2-4x-4)$}
\item\makeblankfill{$3x(5-3x)+4$}
\end{enumerate}
\end{example}
\noindent Note that \makeblank[1.5] are automatically in a factored form.